% =====================================================================
%  ANNEXE C — 14 FICHES DE RÉVISION (une page chacune)
% =====================================================================
\chapter{Quatorze fiches de révision à imprimer}

Ces fiches résument tout ce que le livre développe. Leur usage est simple :
\textbf{imprimez-les} (ou recopiez-les à la main, ce qui est encore plus efficace),
\textbf{affichez-les} là où vous passez du temps --- miroir de la salle de bain, porte du
frigo, mur du bureau, écran de veille du téléphone --- et \textbf{lisez-les à voix haute}
une minute par jour. En trois mois, vous connaîtrez la grammaire anglaise sans jamais
ouvrir un livre de grammaire.

\section*{Fiche 1 --- Les douze temps en une page}
\begin{tableau}[Les temps : forme, emploi, signal]
\footnotesize
\begin{longtable}{@{}p{3.4cm}p{4.6cm}p{3.2cm}p{3.0cm}@{}}
\entete{Temps} & \entete{Emploi} & \entete{Signal} & \entete{Exemple} \tabularnewline
\endfirsthead
\entete{Temps} & \entete{Emploi} & \entete{Signal} & \entete{Exemple} \tabularnewline
\endhead
\hline
\en{present simple} & habitude, vérité générale & \en{every day, always, usually} & \en{I work in Dakar.} \tabularnewline
\en{present continuous} & en cours maintenant, temporaire & \en{now, at the moment, these days} & \en{I'm working now.} \tabularnewline
\en{present perfect} & passé lié au présent, expérience & \en{already, yet, ever, since, for} & \en{I have finished.} \tabularnewline
\en{present perfect cont.} & activité commencée et encore en cours & \en{for, since, all day, how long} & \en{I have been waiting for an hour.} \tabularnewline
\en{past simple} & action terminée à un moment précis & \en{yesterday, last week, in 2019, ago} & \en{I worked yesterday.} \tabularnewline
\en{past continuous} & décor, action interrompue & \en{while, when, at 8 o'clock} & \en{I was working when he called.} \tabularnewline
\en{past perfect} & antériorité dans le passé & \en{before, after, already, by the time} & \en{I had worked before he arrived.} \tabularnewline
\en{past perfect cont.} & durée avant un moment passé & \en{for, since, all morning} & \en{I had been working for hours.} \tabularnewline
\en{future (will)} & décision immédiate, prédiction, promesse & \en{tomorrow, soon, I think} & \en{I'll work tomorrow.} \tabularnewline
\en{future (going to)} & intention, plan, évidence & \en{tonight, next week, look!} & \en{I'm going to work late.} \tabularnewline
\en{future continuous} & en cours à un moment futur & \en{this time tomorrow, at 9} & \en{I'll be working at nine.} \tabularnewline
\en{future perfect} & achevé avant un moment futur & \en{by, by the time, before} & \en{I'll have finished by Friday.} \tabularnewline
\end{longtable}
\end{tableau}

\begin{regle}[La règle d'or des temps]
Trois questions suffisent : \textbf{(1)} l'action est-elle finie et datée ? \ra passé simple.
\textbf{(2)} est-elle finie mais encore liée à maintenant ? \ra \en{present perfect}.
\textbf{(3)} est-elle en train de se dérouler ? \ra un temps en \en{-ing}.
\end{regle}

\section*{Fiche 2 --- Modaux et auxiliaires}
\begin{tableau}[Ce que chaque modal exprime]
\footnotesize
\begin{longtable}{@{}p{3.0cm}p{5.2cm}p{6.0cm}@{}}
\entete{Modal} & \entete{Sens} & \entete{Exemple} \tabularnewline
\endfirsthead
\entete{Modal} & \entete{Sens} & \entete{Exemple} \tabularnewline
\endhead
\hline
\en{can / could} & capacité, possibilité, demande & \en{I can swim. Could you help me?} \tabularnewline
\en{may / might} & permission, probabilité faible & \en{It might rain. May I come in?} \tabularnewline
\en{must} & obligation forte, déduction certaine & \en{You must stop. He must be tired.} \tabularnewline
\en{have to} & obligation extérieure, factuelle & \en{I have to work on Saturdays.} \tabularnewline
\en{should / ought to} & conseil, ce qui serait bien & \en{You should see a doctor.} \tabularnewline
\en{will / would} & futur, volonté, habitude passée, politesse & \en{I'll come. Would you like tea?} \tabularnewline
\en{shall} & proposition, suggestion (surtout \en{I/we}) & \en{Shall we go?} \tabularnewline
\en{used to} & habitude passée révolue & \en{I used to smoke.} \tabularnewline
\en{be used to} & être habitué à (+ \en{-ing}) & \en{I'm used to getting up early.} \tabularnewline
\en{needn't / don't have to} & absence d'obligation & \en{You needn't come.} \tabularnewline
\en{mustn't} & interdiction (\ko pas « pas obligé ») & \en{You mustn't smoke here.} \tabularnewline
\en{can't have / must have} & déduction sur le passé & \en{He can't have left. She must have forgotten.} \tabularnewline
\en{should have} & regret, reproche & \en{You should have told me.} \tabularnewline
\end{longtable}
\end{tableau}

\section*{Fiche 3 --- Poser toutes les questions}
\begin{tableau}[Les six structures interrogatives]
\footnotesize
\begin{longtable}{@{}p{4.4cm}p{9.8cm}@{}}
\entete{Type} & \entete{Structure et exemples} \tabularnewline
\endfirsthead
\entete{Type} & \entete{Structure et exemples} \tabularnewline
\endhead
\hline
\textbf{Oui/non} & Auxiliaire + sujet + verbe. \en{Do you like tea? Are you ready? Can she drive? Have they left? Did he call?} \tabularnewline
\textbf{Mot interrogatif} & Mot en \en{Wh-} + auxiliaire + sujet + verbe. \en{Where do you live? Why did she cry? How long have you been waiting?} \tabularnewline
\textbf{Sujet} (\en{who/what}) & Pas d'auxiliaire \en{do} : \en{Who called you? What happened? Who wants coffee?} \tabularnewline
\textbf{Question tag} & Affirmation + tag négatif (et inversement). \en{You're tired, aren't you? She doesn't smoke, does she? He can swim, can't he?} \tabularnewline
\textbf{Indirecte} & Formule polie + ordre affirmatif. \en{Could you tell me where the station is? Do you know what time it starts?} \tabularnewline
\textbf{Choix} & \en{or} entre deux options. \en{Do you want tea or coffee? Are you coming or staying?} \tabularnewline
\end{longtable}
\end{tableau}

\begin{piege}[Le \en{-s} de la question]
\en{*Does she works?} \ko \ra \en{Does she work?} \ok. L'auxiliaire \en{does} porte déjà la
marque de la 3\textsuperscript{e} personne : le verbe principal reste à la base.
\end{piege}

\section*{Fiche 4 --- Les articles}
\begin{tableau}[Quand utiliser \en{a/an}, \en{the}, ou rien]
\footnotesize
\begin{longtable}{@{}p{3.2cm}p{11.0cm}@{}}
\entete{Cas} & \entete{Règle et exemples} \tabularnewline
\endfirsthead
\entete{Cas} & \entete{Règle et exemples} \tabularnewline
\endhead
\hline
\en{a / an} & Singulier dénombrable, première mention, non identifié. \en{I saw \textbf{a} dog. She's \textbf{an} engineer.} \tabularnewline
\en{the} & Chose identifiée, unique, déjà mentionnée, ou avec superlatif. \en{\textbf{The} dog was black. \textbf{The} sun. \textbf{The} best film.} \tabularnewline
\textbf{rien} & Généralités au pluriel ou indénombrables. \en{\textbf{Dogs} are friendly. I love \textbf{music}. \textbf{Water} is essential.} \tabularnewline
\textbf{rien} & Repas, langues, sports, matières, jours, mois. \en{I had \textbf{breakfast}. She speaks \textbf{English}. I play \textbf{tennis}.} \tabularnewline
\textbf{rien} & Institutions dans leur fonction. \en{go to \textbf{school} / \textbf{bed} / \textbf{church} / \textbf{work} / \textbf{hospital} (BrE)} \tabularnewline
\en{the} & Pays pluriels ou composés, fleuves, mers, déserts, chaînes. \en{\textbf{the} USA, \textbf{the} Netherlands, \textbf{the} Nile, \textbf{the} Sahara, \textbf{the} Alps} \tabularnewline
\textbf{rien} & Pays, villes, continents, lacs, monts isolés. \en{Senegal, Dakar, Africa, Lake Victoria, Mount Kilimanjaro} \tabularnewline
\end{longtable}
\end{tableau}

\section*{Fiche 5 --- Prépositions de temps et de lieu}
\begin{tableau}[\en{at / on / in} --- la règle des trois échelles]
\footnotesize
\begin{longtable}{@{}p{2.2cm}p{4.4cm}p{7.6cm}@{}}
\entete{Prép.} & \entete{Échelle} & \entete{Exemples} \tabularnewline
\endfirsthead
\entete{Prép.} & \entete{Échelle} & \entete{Exemples} \tabularnewline
\endhead
\hline
\en{at} & point précis & \en{at 7 o'clock, at noon, at night, at the weekend (BrE), at Christmas, at the bus stop, at home, at work} \tabularnewline
\en{on} & jour, surface & \en{on Monday, on 3 May, on Monday morning, on my birthday, on the table, on the wall, on the bus, on the plane} \tabularnewline
\en{in} & période longue, contenant & \en{in May, in 2026, in the morning, in summer, in a week, in the box, in the car, in Dakar, in Africa} \tabularnewline
\en{by} & avant ou au plus tard & \en{by Friday, by 6 p.m.} \tabularnewline
\en{until / till} & jusque'à & \en{until midnight, till Friday} \tabularnewline
\en{for / since} & durée / point de départ & \en{for two hours, since 2019} \tabularnewline
\en{during / while} & pendant & \en{during the film, while I was sleeping} \tabularnewline
\en{from\ldots{} to} & de\ldots{} à & \en{from Monday to Friday, from 9 to 5} \tabularnewline
\en{between} & entre deux & \en{between 2019 and 2022, between the two buildings} \tabularnewline
\en{to / into / onto} & mouvement vers & \en{go to school, come into the room, climb onto the roof} \tabularnewline
\en{from / out of / off} & mouvement depuis & \en{from Dakar, out of the car, off the bus} \tabularnewline
\end{longtable}
\end{tableau}

\begin{astuce}[Le moyen mnémotechnique des échelles]
\en{at} = un point (heure, adresse précise). \en{on} = un jour ou une surface.
\en{in} = une boîte (mois, année, ville, pays, contenant). Plus l'échelle est grande,
plus la préposition est « large ».
\end{astuce}

\section*{Fiche 6 --- Comparatifs et superlatifs}
\begin{tableau}[Formation]
\footnotesize
\begin{longtable}{@{}p{3.6cm}p{4.0cm}p{6.6cm}@{}}
\entete{Adjectif} & \entete{Comparatif} & \entete{Superlatif} \tabularnewline
\endfirsthead
\entete{Adjectif} & \entete{Comparatif} & \entete{Superlatif} \tabularnewline
\endhead
\hline
1 syllabe & \en{+ -er} & \en{the + -est} --- \en{tall, taller, the tallest} \tabularnewline
1 syllabe en \en{-e} & \en{+ -r} & \en{the + -st} --- \en{large, larger, the largest} \tabularnewline
consonne + voyelle + consonne & double consonne & \en{big, bigger, the biggest} \tabularnewline
en \en{-y} & \en{-ier} & \en{the -iest} --- \en{happy, happier, the happiest} \tabularnewline
2 syllabes et plus & \en{more +} & \en{the most +} --- \en{careful, more careful, the most careful} \tabularnewline
irréguliers & \en{good/better/best} & \en{bad/worse/worst ; far/further/furthest ; little/less/least ; much/many/more/most} \tabularnewline
égalité & \en{as + adj + as} & \en{She's as tall as me.} \tabularnewline
inégalité & \en{not as + adj + as} & \en{It's not as cold as yesterday.} \tabularnewline
progression & \en{comparative + and + comparative} & \en{It's getting colder and colder.} \tabularnewline
proportion & \en{the + comparative, the + comparative} & \en{The more you practise, the better you get.} \tabularnewline
\end{longtable}
\end{tableau}

\section*{Fiche 7 --- Le discours rapporté}
\begin{tableau}[Le décalage des temps]
\footnotesize
\begin{longtable}{@{}p{5.4cm}p{4.4cm}p{4.4cm}@{}}
\entete{Discours direct} & \entete{Discours rapporté} & \entete{Exemple} \tabularnewline
\endfirsthead
\entete{Direct} & \entete{Rapporté} & \entete{Exemple} \tabularnewline
\endhead
\hline
\en{present simple} & \en{past simple} & \en{``I like it'' \ra She said she liked it.} \tabularnewline
\en{present continuous} & \en{past continuous} & \en{``I'm working'' \ra He said he was working.} \tabularnewline
\en{past simple} & \en{past perfect} & \en{``I saw it'' \ra She said she had seen it.} \tabularnewline
\en{present perfect} & \en{past perfect} & \en{``I've done it'' \ra He said he had done it.} \tabularnewline
\en{will} & \en{would} & \en{``I'll call'' \ra She said she would call.} \tabularnewline
\en{can} & \en{could} & \en{``I can help'' \ra He said he could help.} \tabularnewline
\en{must} & \en{had to} & \en{``I must go'' \ra She said she had to go.} \tabularnewline
\en{may} & \en{might} & \en{``It may rain'' \ra He said it might rain.} \tabularnewline
\en{today} & \en{that day} & --- \tabularnewline
\en{tomorrow} & \en{the next day / the following day} & --- \tabularnewline
\en{yesterday} & \en{the day before / the previous day} & --- \tabularnewline
\en{here} & \en{there} & --- \tabularnewline
\en{this / these} & \en{that / those} & --- \tabularnewline
\en{ago} & \en{before / earlier} & --- \tabularnewline
\end{longtable}
\end{tableau}

\begin{regle}[\en{say} ou \en{tell} ?]
\en{tell} + personne : \en{She \textbf{told me} she was tired.}\\
\en{say} + chose (+ \en{to} personne) : \en{She \textbf{said} she was tired.} ---
\en{She \textbf{said to me}\ldots{}} (plus rare).\\
Jamais \en{*She said me}. Jamais \en{*She told that\ldots{}} sans complément de personne.
\end{regle}

\section*{Fiche 8 --- Les conditionnels}
\begin{tableau}[Les cinq conditionnels]
\footnotesize
\begin{longtable}{@{}p{3.0cm}p{5.0cm}p{6.2cm}@{}}
\entete{Type} & \entete{Structure} & \entete{Emploi et exemple} \tabularnewline
\endfirsthead
\entete{Type} & \entete{Structure} & \entete{Emploi et exemple} \tabularnewline
\endhead
\hline
\textbf{Zéro} & \en{if + present, present} & Vérité générale. \en{If you heat water, it boils.} \tabularnewline
\textbf{1} & \en{if + present, will} & Futur réel ou probable. \en{If it rains, I'll stay home.} \tabularnewline
\textbf{2} & \en{if + past, would} & Imaginaire présent, hypothèse. \en{If I had time, I'd learn Spanish.} \tabularnewline
\textbf{3} & \en{if + past perfect, would have + p.p.} & Regret sur le passé. \en{If I had studied, I would have passed.} \tabularnewline
\textbf{Mixte} & \en{if + past perfect, would + base} & Cause passée, effet présent. \en{If I had taken that job, I'd be rich now.} \tabularnewline
\end{longtable}
\end{tableau}

\begin{piege}[Jamais \en{will} après \en{if}]
\en{*If it will rain\ldots{}} \ko \ra \en{If it rains\ldots{}} \ok.
Exception rare : \en{if} = « si jamais » dans une demande polie ---
\en{If you \textbf{will} wait a moment, I'll check.}
\end{piege}

\section*{Fiche 9 --- Le passif}
\begin{tableau}[Formation à tous les temps]
\footnotesize
\begin{longtable}{@{}p{4.0cm}p{4.4cm}p{5.8cm}@{}}
\entete{Temps} & \entete{Forme passive} & \entete{Exemple} \tabularnewline
\endfirsthead
\entete{Temps} & \entete{Forme passive} & \entete{Exemple} \tabularnewline
\endhead
\hline
\en{present simple} & \en{am/is/are + p.p.} & \en{English is spoken here.} \tabularnewline
\en{present continuous} & \en{am/is/are being + p.p.} & \en{The road is being repaired.} \tabularnewline
\en{past simple} & \en{was/were + p.p.} & \en{The bridge was built in 1990.} \tabularnewline
\en{past continuous} & \en{was/were being + p.p.} & \en{The car was being washed.} \tabularnewline
\en{present perfect} & \en{has/have been + p.p.} & \en{The work has been finished.} \tabularnewline
\en{past perfect} & \en{had been + p.p.} & \en{It had been stolen.} \tabularnewline
\en{future (will)} & \en{will be + p.p.} & \en{The results will be published.} \tabularnewline
\en{future (going to)} & \en{is/are going to be + p.p.} & \en{A school is going to be built.} \tabularnewline
\en{modals} & \en{modal + be + p.p.} & \en{It must be signed. It can be done.} \tabularnewline
\en{modals (past)} & \en{modal + have been + p.p.} & \en{It should have been checked.} \tabularnewline
\end{longtable}
\end{tableau}

\section*{Fiche 10 --- Quantifieurs}
\begin{tableau}[Dénombrable ou indénombrable ?]
\footnotesize
\begin{longtable}{@{}p{3.2cm}p{4.0cm}p{3.6cm}p{3.4cm}@{}}
\entete{Quantifieur} & \entete{S'emploie avec} & \entete{Sens} & \entete{Exemple} \tabularnewline
\endfirsthead
\entete{Quantifieur} & \entete{Avec} & \entete{Sens} & \entete{Exemple} \tabularnewline
\endhead
\hline
\en{many} & dénombrables (surtout nég./interr.) & beaucoup & \en{How many books?} \tabularnewline
\en{much} & indénombrables (surtout nég./interr.) & beaucoup & \en{How much time?} \tabularnewline
\en{a lot of / lots of} & les deux (affirmatif) & beaucoup & \en{a lot of friends / a lot of water} \tabularnewline
\en{(a) few} & dénombrables & quelques / peu & \en{a few days} \tabularnewline
\en{(a) little} & indénombrables & un peu & \en{a little sugar} \tabularnewline
\en{several} & dénombrables & plusieurs & \en{several times} \tabularnewline
\en{some} & les deux (affirmatif) & un peu de, des & \en{some money, some people} \tabularnewline
\en{any} & les deux (nég./interr.) & aucun, du & \en{I don't have any money.} \tabularnewline
\en{enough} & les deux & assez & \en{enough time, enough chairs} \tabularnewline
\en{all / most / no} & les deux & tout, la plupart, aucun & \en{all the students, most people} \tabularnewline
\en{each / every} & dénombrables singuliers & chaque & \en{every morning} \tabularnewline
\en{both / either / neither} & dénombrables (deux) & les deux / l'un ou l'autre / ni l'un ni l'autre & \en{both of us, neither option} \tabularnewline
\end{longtable}
\end{tableau}

\begin{piege}[\en{a few} $\neq$ \en{few}]
\en{a few friends} = quelques amis (positif, j'en ai). \en{few friends} = peu d'amis
(négatif, presque pas). Idem : \en{a little} (un peu) vs \en{little} (presque pas).
\end{piege}

\section*{Fiche 11 --- Prononciation : les dix sons à surveiller}
\begin{tableau}[Les sons qui trahissent un francophone]
\footnotesize
\begin{longtable}{@{}p{2.4cm}p{3.4cm}p{8.4cm}@{}}
\entete{Son} & \entete{Mot modèle} & \entete{Conseil} \tabularnewline
\endfirsthead
\entete{Son} & \entete{Mot modèle} & \entete{Conseil} \tabularnewline
\endhead
\hline
\en{h} aspiré & \en{house, he, hand} & Soufflez vraiment --- le \en{h} anglais existe. \ko \en{*ouse}. \tabularnewline
\en{th} sourd & \en{think, three, thing} & Langue entre les dents, sans voix. Pas \en{sink}. \tabularnewline
\en{th} sonore & \en{this, that, mother} & Même position, avec vibration. Pas \en{dis}. \tabularnewline
\en{w} & \en{we, water, window} & Lèvres arrondies comme pour \en{ou}, puis \en{wé}. Pas \en{v}. \tabularnewline
\en{v} & \en{very, vote} & Dents sur la lèvre --- le vrai \en{v} français. \tabularnewline
\en{r} anglais & \en{red, really} & Langue en arrière, jamais roulée ni grasseyée. \tabularnewline
\en{i} bref & \en{ship, sit, six} & Court et relâché --- différent de \en{sheep, seat}. \tabularnewline
\en{a} de \en{cat} & \en{cat, man, bad} & Bouche ouverte, entre \en{a} et \en{è}. \tabularnewline
\en{ng} final & \en{sing, thing} & \pron{si-ng} --- jamais de \en{g} explosif, jamais de \en{n} simple. \tabularnewline
\en{schwa} & \en{about, teacher} & Le son le plus fréquent de l'anglais : un \en{e} très court et neutre. \tabularnewline
\end{longtable}
\end{tableau}

\begin{regle}[Accent tonique --- trois réflexes]
\textbf{1.} Dans les mots de deux syllabes, l'accent tombe presque toujours sur la
\textbf{première} : \en{TAble, WAter, HAPpy, PREsent} (nom).\\
\textbf{2.} Les suffixes \en{-tion, -sion, -ic, -ity} attirent l'accent sur la syllabe
\textbf{précédente} : \en{naTION, deCIsion, ecoNOmic, uniVERsity}.\\
\textbf{3.} Verbe et nom de même orthographe changent d'accent :
\en{PREsent} (nom) / \en{preSENT} (verbe) ; \en{REcord} / \en{reCORD} ;
\en{OBject} / \en{obJECT}.
\end{regle}

\section*{Fiche 12 --- Orthographe : les 25 mots les plus ratés}
\begin{tableau}[À recopier trois fois chacun]
\footnotesize
\begin{multicols}{2}
\begin{itemize}\setlength{\itemsep}{1pt}
\item \en{because}
\item \en{necessary}
\item \en{separate}
\item \en{definitely}
\item \en{occurred}
\item \en{beginning}
\item \en{believe}
\item \en{receive}
\item \en{foreign}
\item \en{environment}
\item \en{government}
\item \en{different}
\item \en{business}
\item \en{friend}
\item \en{tomorrow}
\item \en{Wednesday}
\item \en{February}
\item \en{restaurant}
\item \en{beautiful}
\item \en{interesting}
\item \en{important}
\item \en{experience}
\item \en{knowledge}
\item \en{language}
\item \en{successful}
\end{itemize}
\end{multicols}
\end{tableau}

\section*{Fiche 13 --- Faux amis éclair (les 30 plus dangereux)}
\begin{tableau}[À relire avant chaque conversation]
\footnotesize
\begin{longtable}{@{}p{3.6cm}p{5.0cm}p{5.6cm}@{}}
\entete{Mot anglais} & \entete{Vrai sens} & \entete{Ce qu'on croit} \tabularnewline
\endfirsthead
\entete{Mot anglais} & \entete{Vrai sens} & \entete{Ce qu'on croit} \tabularnewline
\endhead
\hline
\en{actually} & en fait & actuellement \tabularnewline
\en{eventually} & finalement & éventuellement \tabularnewline
\en{library} & bibliothèque & librairie \tabularnewline
\en{shop} & magasin & shop (informatique) \tabularnewline
\en{coin} & pièce de monnaie & coin (angle) \tabularnewline
\en{chair} & chaise, présidence & chair (viande) \tabularnewline
\en{money} & argent (financier) & monnaie \tabularnewline
\en{currency} & monnaie (système) & courance \tabularnewline
\en{attend} & assister à & attendre \tabularnewline
\en{wait for} & attendre & --- \tabularnewline
\en{afford} & avoir les moyens & affermir \tabularnewline
\en{assist} & aider & assister à \tabularnewline
\en{injury} & blessure & injure \tabularnewline
\en{insult} & insulte & --- (attention : verbe \pron{inSEULT}) \tabularnewline
\en{sensible} & raisonnable & sensible \tabularnewline
\en{sensitive} & sensible & --- \tabularnewline
\en{comprehensive} & complet, exhaustif & compréhensif \tabularnewline
\en{understanding} & compréhensif & --- \tabularnewline
\en{large} & grand (dimensions) & large \tabularnewline
\en{wide} & large & --- \tabularnewline
\en{to travel} & voyager & travailler \tabularnewline
\en{to work} & travailler & --- \tabularnewline
\en{journey} & trajet, voyage & journée \tabularnewline
\en{day} & journée & --- \tabularnewline
\en{car} & voiture & car (autocar) \tabularnewline
\en{coach} & autocar, entraîneur & --- \tabularnewline
\en{to bless} & bénir & blesser \tabularnewline
\en{to hurt / to injure} & blesser & --- \tabularnewline
\en{preservative} & conservateur (alimentaire) & préservatif \tabularnewline
\en{condom} & préservatif & --- \tabularnewline
\end{longtable}
\end{tableau}

\section*{Fiche 14 --- La trousse de secours orale}
\begin{tableau}[Dix phrases qui sauvent une conversation]
\footnotesize
\begin{longtable}{@{}p{7.6cm}p{6.6cm}@{}}
\entete{Phrase} & \entete{À quoi elle sert} \tabularnewline
\endfirsthead
\entete{Phrase} & \entete{À quoi elle sert} \tabularnewline
\endhead
\hline
\en{Sorry, could you say that again?} & Gagner une seconde chance de comprendre. \tabularnewline
\en{Could you speak more slowly, please?} & Ralentir l'interlocuteur sans l'offenser. \tabularnewline
\en{How do you say ``\ldots{}'' in English?} & Combler un mot manquant. \tabularnewline
\en{You mean the thing you use to\ldots{} ?} & Décrire un objet dont on ignore le nom. \tabularnewline
\en{Let me put it another way.} & Reformuler au lieu de se bloquer. \tabularnewline
\en{What I'm trying to say is\ldots{}} & Rattraper une phrase mal partie. \tabularnewline
\en{Just a second --- let me think.} & Acheter du temps légitimement. \tabularnewline
\en{I'm not sure I follow you.} & Signaler une incompréhension sans s'excuser. \tabularnewline
\en{Could you give me an example?} & Transformer un mot abstrait en concret. \tabularnewline
\en{My English isn't perfect yet, but I'm working on it.} & Désamorcer la pression, en une phrase. \tabularnewline
\end{longtable}
\end{tableau}

\begin{mission}[Annexe C --- le rituel des fiches]
Imprimez les quatorze fiches. Collez-en \textbf{deux} sur le mur de votre espace de travail.
Chaque matin, lisez à voix haute les deux fiches affichées : une minute trente suffit.
Tous les sept jours, remplacez-les par les deux suivantes. En sept semaines, les quatorze
fiches sont passées sept fois devant vos yeux --- et la grammaire anglaise cesse d'être une
liste de règles pour devenir un automatisme.
\end{mission}
